
\subsubsection{5.1 Foundation: Runtime Error Prevalence (H-E1)}

\begin{table}[htbp]
\centering
\resizebox{\columnwidth}{!}{%
\begin{tabular}{lll}
\toprule
Category & Count & Percentage \\
\midrule
Runtime Error & 304 & 60.8\% \\
Syntax Error & 193 & 38.6\% \\
Wrong Output & 3 & 0.6\% \\
Timeout & 0 & 0\% \\
**Total Failures** & **500** & **100\%** \\
\bottomrule
\end{tabular}
}
\end{table}

Runtime error prevalence was 60.8\% (95\% CI: 56.5\%-65.0\%), exceeding the 30\% threshold specified in H-E1. The gate condition was satisfied.

Note: The 0\% pass rate for initial code generation indicates that CodeLlama-7B-Instruct did not successfully solve any MBPP problems in this configuration. This does not affect the granularity comparison, which focuses on the distribution of error types among failures.

\subsubsection{5.2 Main Effect: Granularity and Repair Success (H-M1)}

\begin{table}[htbp]
\centering
\resizebox{\columnwidth}{!}{%
\begin{tabular}{llll}
\toprule
Level & Successes & Rate & 95\% CI \\
\midrule
G0 & 127/304 & 41.8\% & [36.3\%, 47.4\%] \\
G1 & 124/304 & 40.8\% & [35.4\%, 46.4\%] \\
G2 & 56/304 & 18.4\% & [14.4\%, 23.2\%] \\
G3 & 51/304 & 16.8\% & [12.9\%, 21.4\%] \\
G4 & 69/304 & 22.7\% & [18.3\%, 27.8\%] \\
\bottomrule
\end{tabular}
}
\end{table}

\textbf{ANOVA Results:}

\begin{table}[htbp]
\centering
\resizebox{\columnwidth}{!}{%
\begin{tabular}{ll}
\toprule
Statistic & Value \\
\midrule
F-statistic & 23.89 \\
p-value & 3.5e-19 \\
Eta-squared & 0.059 \\
\bottomrule
\end{tabular}
}
\end{table}

The ANOVA indicates a statistically significant effect of granularity on repair success. The eta-squared of 0.059 corresponds to a medium effect size by conventional criteria.

\subsubsection{5.3 Post-Hoc Pairwise Comparisons (Tukey HSD)}

\begin{table}[htbp]
\centering
\resizebox{\columnwidth}{!}{%
\begin{tabular}{llll}
\toprule
Comparison & Difference & p-value & Significant \\
\midrule
G0 vs G1 & +1.0pp & 0.999 & No \\
G0 vs G2 & +23.4pp & 5.9e-10 & Yes \\
G0 vs G3 & +25.0pp & 2.5e-11 & Yes \\
G0 vs G4 & +19.1pp & 8.3e-07 & Yes \\
G1 vs G2 & +22.4pp & 3.5e-09 & Yes \\
G1 vs G3 & +24.0pp & 1.7e-10 & Yes \\
G1 vs G4 & +18.1pp & 3.6e-06 & Yes \\
G2 vs G3 & +1.6pp & 0.990 & No \\
G2 vs G4 & -4.3pp & 0.747 & No \\
G3 vs G4 & -5.9pp & 0.452 & No \\
\bottomrule
\end{tabular}
}
\end{table}

The results indicate two clusters: G0 and G1 form a high-success group (~41\%), while G2, G3, and G4 form a lower-success group (~17-23\%). Within-cluster differences are not statistically significant.

\subsubsection{5.4 Directional Test: G0 vs G3 (H-M2)}

H-M2 predicted that G3 would outperform G0 by at least 10 percentage points.

\begin{table}[htbp]
\centering
\resizebox{\columnwidth}{!}{%
\begin{tabular}{ll}
\toprule
Metric & Value \\
\midrule
G0 success rate & 41.8\% \\
G3 success rate & 16.8\% \\
Difference (G3 - G0) & -25.0pp \\
95\% CI & [-32.0pp, -18.0pp] \\
McNemar chi-squared & 77 \\
p-value & 5.23e-22 \\
\bottomrule
\end{tabular}
}
\end{table}

\textbf{Contingency table (paired outcomes):}

\begin{table}[htbp]
\centering
\resizebox{\columnwidth}{!}{%
\begin{tabular}{lll}
\toprule
 & G3 Success & G3 Failure \\
\midrule
G0 Success & 50 & 77 \\
G0 Failure & 1 & 176 \\
\bottomrule
\end{tabular}
}
\end{table}

The difference is in the opposite direction to the hypothesis. G0 outperforms G3 by 25 percentage points. Of the 78 discordant pairs, 77 favor G0 while only 1 favors G3. H-M2 is not supported.

\subsubsection{5.5 Non-Monotonicity Test: G3 vs G4 (H-M3)}

H-M3 predicted that G4 would not significantly outperform G3 (difference within 2\%).

\begin{table}[htbp]
\centering
\resizebox{\columnwidth}{!}{%
\begin{tabular}{ll}
\toprule
Metric & Value \\
\midrule
G3 success rate & 16.8\% \\
G4 success rate & 22.7\% \\
Difference (G4 - G3) & +5.9pp \\
McNemar p-value & 4.0e-05 \\
\bottomrule
\end{tabular}
}
\end{table}

\textbf{Contingency table:}

\begin{table}[htbp]
\centering
\resizebox{\columnwidth}{!}{%
\begin{tabular}{lll}
\toprule
 & G4 Success & G4 Failure \\
\midrule
G3 Success & 50 & 1 \\
G3 Failure & 19 & 234 \\
\bottomrule
\end{tabular}
}
\end{table}

G4 significantly outperforms G3 (p < 0.001), with the difference (5.9pp) exceeding the 2\% margin specified in the hypothesis. H-M3 is not supported.

\subsubsection{5.6 Summary of Hypothesis Outcomes}

\begin{table}[htbp]
\centering
\resizebox{\columnwidth}{!}{%
\begin{tabular}{llll}
\toprule
Hypothesis & Prediction & Result & Status \\
\midrule
H-E1 & Runtime errors >= 30\% & 60.8\% & Supported \\
H-M1 & ANOVA p < 0.05 & p = 3.5e-19 & Supported \\
H-M2 & G3 >= G0 + 10pp & G3 = G0 - 25pp & Not supported \\
H-M3 & G4 <= G3 + 2\% & G4 = G3 + 5.9\% & Not supported \\
\bottomrule
\end{tabular}
}
\end{table}
